\section{Related Work}
\label{sec:related}

Mature production frameworks have driven large-scale MoE training.
\megatron~\cite{shoeybi2020megatronlmtrainingmultibillionparameter,narayanan2021efficientlargescalelanguagemodel}
established the pipeline-parallel Transformer recipe most
subsequent frameworks build on, and
\deepspeed~\cite{deepspeed,rajbhandari2020zeromemoryoptimizationstraining}
introduced ZeRO sharded-optimizer techniques. Both rely on
layered designs with plugin systems, registry-based indirection,
and compiled extension chains, delivering broad model coverage
and peak throughput on production hardware.
\torchtitan~\cite{liang2025torchtitanonestoppytorchnative} more
recently scaled PyTorch-native training to multi-thousand-GPU
clusters, sharing \sys's Python-native goal but without making
agent-native design a primary axis or matching production-framework
throughput on our configurations. \sys takes a different point on
this tradeoff curve with agent-native design as a first-class
goal, achieving higher agent-task efficiency along Agent Turns,
Output Tokens, etc., while matching the training throughput of
production frameworks.

AI coding agents combine a tool-using language model with a small,
stable vocabulary of primitives (read, edit, shell, search).
Recent work designs better agent flows and harnesses for software
engineering on top of these
primitives~\cite{yao2023react,yang2024sweagentagentcomputerinterfacesenable,zhang2024autocoderoverautonomousprogramimprovement}.
\sys takes a complementary direction: rather than improving the
agent flow, we design the software framework so that existing
flows work better, lowering agent cost without changing the agent.

Existing agent benchmarks score capability on fixed codebases and
tasks: SWE-bench~\cite{jimenez2024swebench} on GitHub-issue
resolution, MLE-bench~\cite{chan2025mlebench} on Kaggle-style ML
engineering, and
HumanEval~\cite{chen2021evaluatinglargelanguagemodels} on
function-level code generation. Closer to ML systems engineering,
FlashInfer-Bench~\cite{xing2026flashinferbenchbuildingvirtuouscycle}
and KernelBench~\cite{ouyang2025kernelbench} target inference
operators and GPU kernels respectively. The agent is the variable;
aggregate task correctness is the metric. \sysbench inverts this: holding
the agent and task fixed, we vary the training framework so
that differences in agent cost and task outcome isolate framework
design. This axis is complementary to
existing capability benchmarks.
